\begin{abstract}
Parallel software agents increasingly edit one body of files from different
processes and machines.  They need Git-like retained history and conflict
visibility, but synchronization must be continuous rather than organized around
explicit repository commits.  Versioning, synchronization, encrypted remote
storage, and portable execution are commonly separate layers whose boundaries
give one update several identities: a filesystem block, a version-control
delta, a synchronization object, and a cloud object.  We present \sfs{}, a
filesystem substrate motivated by this \emph{agentic-time} setting that instead
makes a position-addressed \emph{fragment version map} the common
representation.  Each unit carries a
sparse version vector and, for every fragment position, the causal dot of the
bytes currently selected there.  The map is simultaneously a retained local
lineage, a synchronization manifest, and the basis for sub-file conflict
classification.  Concurrent updates to disjoint fragments merge without a
content-specific CRDT.  Concurrent updates to the same fragment remain visible
as separate strains.

Around this representation, \sfs{} combines deterministic fixed-size chunking,
atomic root publication, retained versions and named cuts, client-side-encrypted
synchronization, and an epoch-versioned multi-writer authorization model.  The
same synchronization protocol operates against either a blind cloud store or a
live peer.  In the secure deployment profile, content and metadata are
client-side AES-GCM encrypted and accepted record frontiers are authorized by
an owner-signed writer set.  The service receives neither plaintext nor the keys
needed to derive it.  A single v12 container format is implemented by a
memory-safe Rust core, a FUSE adapter over host-file containers, an independent
in-kernel C driver over raw devices, a C ABI, and experimental in-memory
WebAssembly bindings.  The same logical state therefore crosses an execution
boundary (kernel/userspace/browser) and a trust boundary (local media/encrypted
service) without conversion into a second storage model.

Deterministic multi-replica scenarios exercise disjoint-fragment convergence,
same-fragment strain preservation and resolution, third-replica discovery, and
bidirectional synchronization with a live TLS peer.  They establish the
mechanisms in small controlled topologies, not behavior at agent-fleet scale.
We separately evaluate the mounted Linux paths over 14 I/O sizes using ten
fault-checked runs per cell.  Under the reported single-threaded protocol,
kernel-path geometric means are $1.22\times$ ext4 for plaintext writes,
$0.98\times$ for plaintext cold reads, $0.98$--$0.99\times$ ext4-on-LUKS for
encrypted writes,
and $1.11$--$1.20\times$ for encrypted cold reads.  The FUSE path reaches
$3.11\times$ \texttt{fuse2fs} for plaintext writes and $1.70\times$ for
plaintext cold reads, but only $0.33$--$0.38\times$ \texttt{gocryptfs} for
encrypted cold reads.  These results establish a competitive native path and
identify the userspace encrypted-read path as the principal performance limit.
They do not constitute a general workload or scalability evaluation.  As a
secondary result, an annex reports a NoSQL projection built without changing
the core representation.
\end{abstract}
